\documentclass[10pt,a4paper,titlepage]{jreport} % ドキュメントクラスを1つに統一


\usepackage[top=25truemm,bottom=25truemm,left=20truemm,right=20truemm]{geometry} % 余白設定
\usepackage[dvipdfmx]{graphicx} % 画像の挿入
\usepackage{graphicx}
\usepackage{listings} % コードリスト用
\usepackage{jlisting} % 日本語対応のリスト用
\usepackage{jverb} % 日本語のベタ打ち
\usepackage{booktabs}
\usepackage{pgfplots} % グラフ描画用パッケージ
\pgfplotsset{compat=1.18} % バージョン指定
\usepackage{float}
\parindent = 0pt % 段落の字下げを無効化

\usepackage{titlesec}
\usepackage{etoolbox}
\makeatletter
\patchcmd{\chapter}{\if@openleft\cleardoublepage\else\if@openright\cleardoublepage\else\clearpage\fi\fi}{}{}{}
\makeatother
\lstset{breaklines=true, postbreak=\mbox{$\hookrightarrow$}\space, keepspaces=true, escapeinside={\%*}{*)}}

\titleformat{\chapter}[hang] % 章のフォーマットを変更
  {\normalfont\huge\bfseries} % フォント設定
  {\thechapter} % 番号部分の表示形式
  {1em} % 番号とタイトルの間のスペース
  {} % タイトルの前に入れる内容（ここでは空）
\usepackage{xcolor}
\usepackage{amsmath}
\usepackage{amssymb} % ←これで \mathbb{R} が使える

\lstset{
  language=Matlab,              % 言語はMATLAB
  basicstyle=\ttfamily\small,   % フォントは等幅、サイズ小さめ
  keywordstyle=\color{blue},    % キーワード（function, endとか）を青
  commentstyle=\color{green!50!black}, % コメントを深緑
  stringstyle=\color{red!70!black},    % 文字列（'...'）を赤系
  numbers=left,                 % 行番号を左に表示
  numberstyle=\tiny\color{gray},% 行番号をグレーで小さく
  stepnumber=1,                 % 毎行行番号つける
  numbersep=10pt,               % コードとの間隔
  backgroundcolor=\color{gray!10}, % 背景うっすらグレー
  frame=single,                 % 枠をつける
  rulecolor=\color{black},      % 枠線は黒
  breaklines=true,              % 長い行を折り返す
  breakatwhitespace=false,      % スペースでだけ折り返すか（falseにして自然な折り返し）
  showspaces=false,             % スペースは特別に表示しない
  showstringspaces=false,       % 文字列中のスペースも普通に表示
  tabsize=2                     % タブ幅2
}


\title{数値計算：課題3} % タイトル
\author{
  学生番号：242C2016  氏名：奥村直 \\
  \\
  知的システム工学科システム制御コース
  } % 著者
\date{\today} % 日付

\begin{document}

\maketitle

\section*{課題}

\subsection*{問題設定}
ガウスの単純消去法を用いて，以下の連立一次方程式を解く
プログラムを作成する。ただし，$n$ は学生番号の下1桁とし，
$n=6$ の場合を考える。

\[
\begin{bmatrix}
n & 2 & 3 \\
2 & 2n & 1 \\
3 & 1 & 3n
\end{bmatrix}
\begin{bmatrix}
x_1 \\ x_2 \\ x_3
\end{bmatrix}
=
\begin{bmatrix}
1 \\ 0 \\ 0
\end{bmatrix}
\]

$n=1$ を代入すると，
\[
\begin{bmatrix}
1 & 2 & 3 \\
2 & 2 & 1 \\
3 & 1 & 3
\end{bmatrix}
\begin{bmatrix}
x_1 \\ x_2 \\ x_3
\end{bmatrix}
=
\begin{bmatrix}
1 \\ 0 \\ 0
\end{bmatrix}
\]
となる。

\subsection*{ガウスの単純消去法}
まず前進消去により係数行列を上三角行列に変換し，
その後，後退代入によって解を求める。

前進消去では，$k$ 行目を基準として，
\[
m_i = \frac{a_{ik}}{a_{kk}}
\]
を用い，各行から基準行の倍を引くことで，
$a_{ik}$ を $0$ にする。

\subsection*{使用したプログラム}
以下に，本課題で使用した C 言語のプログラムを示す。

\begin{lstlisting}
#include <stdio.h>

int main(void) {
    int n = 3;
    double a[3][4] = {
        {6, 2, 3, 1},
        {2, 12, 1, 0},
        {3, 1, 18, 0}
    };
    double x[3];

    // 前進消去
    for (int k = 0; k < n - 1; k++) {
        for (int i = k + 1; i < n; i++) {
            double m = a[i][k] / a[k][k];
            for (int j = k; j <= n; j++) {
                a[i][j] -= m * a[k][j];
            }
        }
    }

    // 後退代入
    for (int i = n - 1; i >= 0; i--) {
        x[i] = a[i][n];
        for (int j = i + 1; j < n; j++) {
            x[i] -= a[i][j] * x[j];
        }
        x[i] /= a[i][i];
    }

    // 結果出力
    for (int i = 0; i < n; i++) {
        printf("x_%d = %.6f\n", i + 1, x[i]);
    }

    return 0;
}
\end{lstlisting}

\subsection*{計算結果}
プログラムを実行した結果，以下の解が得られた。
\[
x_1 = 0.191622
x_2 = -0.029412
x_3 = -0.030303
\]

\end{document}